% UTF-8 / upLaTeX source.
% Compile example:
%   uplatex fmo_matrix_projection_kaisetsu.tex
%   dvipdfmx fmo_matrix_projection_kaisetsu.dvi

\documentclass[uplatex,dvipdfmx,11pt]{jsarticle}

\usepackage{amsmath,amssymb}
\usepackage{geometry}
\usepackage{array}
\usepackage{longtable}

\geometry{a4paper,margin=25mm}

\newcommand{\Hcore}{\mathbf H}
\newcommand{\Venv}{\mathbf V}
\newcommand{\Smat}{\mathbf S}
\newcommand{\Amat}{\mathbf A}
\newcommand{\Bmat}{\mathbf B}
\newcommand{\Tmat}{\mathbf T}
\newcommand{\Cmat}{\mathbf C}
\newcommand{\Cr}{\mathbf C_{\mathrm r}}
\newcommand{\Fmat}{\mathbf F}
\newcommand{\Fr}{\mathbf F_{\mathrm r}}
\newcommand{\Sr}{\mathbf S_{\mathrm r}}
\newcommand{\Dmat}{\mathbf D}
\newcommand{\Dr}{\mathbf D_{\mathrm r}}
\newcommand{\Gmat}{\mathbf G}

\title{FMO境界処理と矩形変換による混成軌道の除去}
\author{}
\date{作成日: 2026年7月5日}

\begin{document}
\maketitle

\section*{この文書の目的}

FMO法で分子をフラグメントに分けたとき、切断した結合の方向を向いた
「使ってはいけない混成軌道」を、どうやって計算から取り除くかを記述する。

結論を先に言うと、従来のように大きなレベルシフト \(B\) でその軌道を無理に押し上げる
必要はない。代わりに、最初からその方向を含まない小さい空間でSCFを解けばよい。
この操作を、ここでは「矩形変換によるdeflation」と呼ぶ。

\section{まず何が問題なのか}

FMOはFragment Molecular Orbitalの略で、大きな分子をいくつかの小さなフラグメントに
分けて量子化学計算を行う方法である。タンパク質や高分子のように、全体を一度に扱う
のが重い系では、この考え方が重要になる。

ただし、分子をフラグメントに分けるとき、共有結合を切ることがある。たとえば、炭素と
炭素の単結合を途中で切って、左側をフラグメント \(x\)、右側を別のフラグメントにする
状況である。

共有結合を切ると、切断点には「もともと結合を作っていた方向」が残る。炭素が
\(sp^3\) 混成している場合なら、4本の \(sp^3\) 混成軌道のうち、切った結合の方向を
向いた1本が問題になる。

模式的には次のような状況である。

\begin{verbatim}
Fragment x side                 Other fragment side

   allowed directions -- C_alpha ---- [ cut-direction sp3 hybrid ]
                                           ^
                                  not allowed in x calculation
\end{verbatim}

この「切断方向の混成軌道」を、この文書では
\[
  \varphi_i^h
\]
と書く。複数本ある場合は \(i=1,2,\ldots\) と番号を付ける。

フラグメント \(x\) の計算でこの軌道をそのまま使える状態にしておくと、フラグメント
\(x\) の電子が、本来は相手フラグメント側に属する方向へ入り込んでしまう。これは
物理的に望ましくない。

したがって、FMOの境界処理では次のことを行う必要がある。

\begin{quote}
切断した結合の方向を向いた混成軌道を、フラグメントのSCF変分空間から取り除く。
\end{quote}

ここで「変分空間」とは、SCFが分子軌道を探すときに動かしてよい自由度の集合である。
そこに不要な方向が残っていると、SCFはその方向も使ってしまう可能性がある。

\section{最低限の量子化学の言葉}

\subsection{AOとMO}

量子化学計算では、分子軌道MOを原子軌道AOの線形結合として表す。
\[
  \psi_p = \sum_\mu C_{\mu p}\chi_\mu .
\]
ここで、\(\chi_\mu\) はAO、\(\psi_p\) はMO、\(C_{\mu p}\) はAOをどれだけ混ぜて
MOを作るかを表す係数である。

すべてのMO係数をまとめた行列を \(\Cmat\) と書く。SCFは、ざっくり言えば、この
\(\Cmat\) を調整して電子のエネルギーを下げる手続きである。

\subsection{AOは普通、直交していない}

AO同士の重なりは
\[
  S_{\mu\nu} = \langle \chi_\mu | \chi_\nu \rangle
\]
で表す。これを集めたものが重なり行列 \(\Smat\) である。

もしAOが完全に直交していれば \(\Smat=\mathbf I\) である。しかし実際の原子中心基底
では、異なる原子上のAOも空間的に重なる。そのため、一般には
\[
  \Smat \ne \mathbf I
\]
である。この「AOが非直交である」という事実が、後で式に \(\Smat\) が何度も出てくる
理由である。

\subsection{除去したい混成軌道}

切断方向の混成軌道もAOの線形結合で表せる。
\[
  \varphi_i^h = \sum_\mu A_{\mu i}\chi_\mu .
\]
この係数を列に並べた行列を \(\Amat\) と書く。AOの数を \(N\)、除去したい混成軌道の
本数を \(m\) とすると、\(\Amat\) は \(N\times m\) の行列である。たとえば、
切断方向の \(sp^3\) 混成軌道が1本だけなら \(m=1\) で、\(\Amat\) は1列のベクトルである。

\section{「使ってはいけない方向を使っていない」とはどういう意味か}

あるMO \(\psi_p\) が、除去したい混成軌道 \(\varphi_i^h\) を含んでいるかどうかは、
両者の内積で見ればよい。
\[
  \langle \varphi_i^h | \psi_p \rangle .
\]
これがゼロなら、\(\psi_p\) は \(\varphi_i^h\) 方向の成分を持っていない。

すべての除去対象 \(i\) とすべてのMO \(p\) についてゼロにしたいので、行列で書くと
\[
  \Amat^\dagger \Smat \Cmat = \mathbf 0
\]
である。

ここに \(\Smat\) が入る理由は、AOが非直交だからである。実際、
\begin{align*}
\langle \varphi_i^h | \psi_p \rangle
&=
\left\langle
\sum_\mu A_{\mu i}\chi_\mu
\middle|
\sum_\nu C_{\nu p}\chi_\nu
\right\rangle \\
&=
\sum_{\mu\nu} A_{\mu i}^{*}
\langle \chi_\mu | \chi_\nu \rangle
C_{\nu p} \\
&=
\sum_{\mu\nu} A_{\mu i}^{*} S_{\mu\nu} C_{\nu p}.
\end{align*}
これを行列でまとめると \(\Amat^\dagger\Smat\Cmat\) になる。

したがって、境界処理の数学的な目的は次の一行にまとめられる。
\[
  \boxed{\Amat^\dagger \Smat \Cmat = \mathbf 0}
\]
これは「どのMOも、切断方向の混成軌道を含まない」という意味である。

\section{従来法: レベルシフトで押し上げる}

従来の方法では、使ってはいけない軌道に非常に大きなエネルギーを与える。
考え方は単純である。

\begin{quote}
その方向を使うとエネルギー的に大損するようにして、SCFがその方向を選ばないようにする。
\end{quote}

この方法をレベルシフト射影、またはhybrid orbital projectionと呼ぶ。

コアハミルトニアンを \(\Hcore\)、環境静電ポテンシャルを \(\Venv\) とすると、
レベルシフト込みの行列は次の形になる。
\[
  \widetilde{\Hcore}
  =
  \Hcore+\Venv+\Smat\Amat\Bmat\Amat^\dagger\Smat .
\]
ここで \(\Bmat\) は大きな正の値を並べた対角行列である。
\[
  \Bmat = \mathrm{diag}(B_1,B_2,\ldots,B_m).
\]
典型的には、\(B_i\) として非常に大きな値を入れる。これにより、
\(\varphi_i^h\) 方向の軌道エネルギーが大きく上がり、被占有軌道から追い出されることを
期待する。

\subsection{\(\Smat\Amat\Bmat\Amat^\dagger\Smat\) という形になる理由}

射影項の成分表示は
\[
  P_{\mu\nu}
  =
  \sum_i B_i
  \langle \chi_\mu | \varphi_i^h \rangle
  \langle \varphi_i^h | \chi_\nu \rangle
\]
である。

混成軌道 \(\varphi_i^h\) はAO係数 \(\Amat\) で表されるが、AOは非直交である。そのため、
\[
  \langle \chi_\mu | \varphi_i^h \rangle
  =
  (\Smat\Amat)_{\mu i}
\]
となる。同様に右側の内積からも \(\Smat\) が出る。したがって、
\[
  \mathbf P = \Smat\Amat\Bmat\Amat^\dagger\Smat
\]
となる。

もしAOが最初から直交していて \(\Smat=\mathbf I\) なら、射影項は単に
\[
  \Amat\Bmat\Amat^\dagger
\]
になる。つまり、左右に出てくる2つの \(\Smat\) は、非直交AOを使っていることの結果である。

\section{レベルシフト法の弱点}

\(B_i\) を大きくすれば、不要な混成軌道の成分は小さくなる。しかし、有限の \(B_i\) では
完全にはゼロにならない。残る成分は大まかには
\[
  O(B^{-1})
\]
の大きさである。

一方で、\(B_i\) を大きくしすぎると、行列の中に極端に大きな値が入る。すると条件数が
悪化し、数値的に不安定になる。SCFの収束にも悪影響が出る可能性がある。

\begin{center}
\begin{tabular}{|p{0.22\linewidth}|p{0.34\linewidth}|p{0.34\linewidth}|}
\hline
\textbf{\(B\) の選び方} & \textbf{良い点} & \textbf{悪い点} \\
\hline
小さめ & 数値的には比較的安定 & 不要成分が残りやすい \\
\hline
大きめ & 不要成分は小さくなる & 条件数が悪化しやすい \\
\hline
無限大 & 理想的には完全除去 & 実際の計算では使えない \\
\hline
\end{tabular}
\end{center}

この問題を避けるには、そもそも \(B\) を使わない方法が必要である。

\section{新しい考え方: 罰するのではなく、最初から選べなくする}

レベルシフト法は「使ってはいけない方向を選んだら罰する」方法である。一方、矩形変換
による方法では、発想を変える。

\begin{quote}
使ってはいけない方向を、変数の中から最初から消しておく。
\end{quote}

SCFが動かす変数はMO係数 \(\Cmat\) である。そこで、\(\Cmat\) を自由に直接動かすのでは
なく、より小さい係数行列 \(\Cr\) を使って
\[
  \Cmat = \Tmat\Cr
\]
と表す。

ここで \(\Tmat\) は \(N\times(N-m)\) の縦長の行列である。列数が \(N\) ではなく
\(N-m\) なので、自由度が \(m\) 個減っている。このため「矩形変換」と呼ぶ。

\(\Tmat\) は、次の条件を満たすように作る。
\[
  \boxed{\Amat^\dagger\Smat\Tmat=\mathbf 0}
\]
すると、
\begin{align*}
  \Amat^\dagger\Smat\Cmat
  &=
  \Amat^\dagger\Smat\Tmat\Cr \\
  &=
  \mathbf 0\,\Cr \\
  &=
  \mathbf 0 .
\end{align*}
これは非常に重要である。\(\Cr\) がどんな値になっても、できあがる
\(\Cmat=\Tmat\Cr\) は必ず
\[
  \Amat^\dagger\Smat\Cmat=\mathbf 0
\]
を満たす。つまり、SCFがどのように進んでも、不要な混成軌道は最初から使えない。

\section{\(\Tmat\) は何を表しているのか}

\(\Tmat\) は、許された方向だけを並べた行列である。除去したい方向を表す条件は
\[
  \Amat^\dagger\Smat\mathbf v=\mathbf 0
\]
である。これを満たすベクトル \(\mathbf v\) は、除去したい混成軌道と \(\Smat\) の意味で
直交している。

このような \(\mathbf v\) 全体の空間を
\[
  \ker(\Amat^\dagger\Smat)
\]
と書く。\(\ker\) はkernel、つまり核または零空間である。

\(\Tmat\) の列は、この核を張る基底である。言い換えると、\(\Tmat\) は「不要な混成軌道と
直交する方向だけを並べた行列」である。

AO空間全体の次元が \(N\)、取り除く混成軌道が \(m\) 本なら、残る許容空間の次元は
\[
  N-m
\]
である。そのため、\(\Tmat\) は \(N\times(N-m)\) の行列になる。

\section{直感的なたとえ}

2つの方法の違いは、次のように考えると分かりやすい。

\begin{center}
\begin{tabular}{|p{0.22\linewidth}|p{0.35\linewidth}|p{0.33\linewidth}|}
\hline
\textbf{方法} & \textbf{たとえ} & \textbf{何が起きるか} \\
\hline
レベルシフト & 入ってはいけない道に高い坂を作る & 坂が有限なら、少しは入り込む可能性がある \\
\hline
矩形変換 & 地図からその道を消す & そもそもその方向へ進めない \\
\hline
\end{tabular}
\end{center}

レベルシフトは「罰する」方法である。矩形変換は「選択肢から消す」方法である。この違いに
より、矩形変換では不要成分が原理的にゼロになる。

\section{矩形変換を使うとレベルシフト項は消える}

レベルシフト込みの行列は
\[
  \widetilde{\Hcore}
  =
  \Hcore+\Venv+\Smat\Amat\Bmat\Amat^\dagger\Smat
\]
であった。これを許された空間へ変換する。つまり、左から \(\Tmat^\dagger\)、右から
\(\Tmat\) を掛ける。
\begin{align*}
\Tmat^\dagger\widetilde{\Hcore}\Tmat
&=
\Tmat^\dagger(\Hcore+\Venv)\Tmat \\
&\quad+
\Tmat^\dagger\Smat\Amat\Bmat\Amat^\dagger\Smat\Tmat .
\end{align*}
第2項は次のように書き換えられる。
\begin{align*}
&\Tmat^\dagger\Smat\Amat\Bmat\Amat^\dagger\Smat\Tmat \\
&\quad=
(\Amat^\dagger\Smat\Tmat)^\dagger
\Bmat
(\Amat^\dagger\Smat\Tmat).
\end{align*}
しかし、\(\Tmat\) は
\[
  \Amat^\dagger\Smat\Tmat=\mathbf 0
\]
を満たすように作っている。したがって、
\[
  (\Amat^\dagger\Smat\Tmat)^\dagger
  \Bmat
  (\Amat^\dagger\Smat\Tmat)
  =
  \mathbf 0
\]
である。

よって、縮約空間で使う行列は
\[
  \boxed{
  \mathbf H_{\mathrm r}
  =
  \Tmat^\dagger(\Hcore+\Venv)\Tmat
  }
\]
になる。ここには、もう \(B\) が出てこない。

これは偶然ではない。\(B\) が担当していた「不要な方向を排除する」という役割を、
\(\Tmat\) の条件
\[
  \Amat^\dagger\Smat\Tmat=\mathbf 0
\]
が完全に引き受けたからである。

\section{縮約空間でのSCF方程式}

通常のRoothaan--Hall方程式は、AOが非直交なので一般化固有値問題になる。矩形変換後は、
縮約空間で次の式を解く。
\[
  \Fr\Cr=\Sr\Cr\boldsymbol\varepsilon .
\]
ここで、
\[
  \Sr=\Tmat^\dagger\Smat\Tmat
\]
であり、
\[
  \Fr
  =
  \Tmat^\dagger
  \left[
    \Hcore+\Venv+\Gmat[\Dmat]
  \right]
  \Tmat
\]
である。\(\Gmat[\Dmat]\) は二電子項、つまりCoulomb項とexchange項を表す。

密度行列も縮約空間から元のAO空間へ戻して使う。
\[
  \Dmat = \Tmat\Dr\Tmat^\dagger .
\]
このようにしておけば、二電子積分やFock行列の構築は元のAO空間の量として扱いつつ、
SCFの自由度だけを \(N-m\) 次元に減らせる。

\section{\(S\)-正規直交化}

\(\Tmat\) の列は、除去したい軌道と \(\Smat\) の意味で直交していれば十分である。ただし、
そのままだと縮約空間の重なり行列
\[
  \Sr=\Tmat^\dagger\Smat\Tmat
\]
が単位行列にならないことがある。

そこで、\(\Tmat\) の列を \(\Smat\) の計量で正規直交化して、
\[
  \Tmat^\dagger\Smat\Tmat=\mathbf I
\]
となるようにできる。この場合、縮約空間のSCF方程式は
\[
  \Fr\Cr=\Cr\boldsymbol\varepsilon
\]
という普通の固有値問題になる。

実装では、まず
\[
  \ker(\Amat^\dagger\Smat)
\]
の基底を求め、それに対してLoewdin直交化や正準直交化を行って \(\Tmat\) を作る。

\section{\(sp^3\) 炭素での具体例}

炭素の価電子AOを単純化して
\[
  s,\quad p_x,\quad p_y,\quad p_z
\]
の4つだけとする。この4つから、4本の \(sp^3\) 混成軌道を作れる。

理想的な場合、たとえば次のように書ける。
\begin{align*}
h_1 &= \frac12(s+p_x+p_y+p_z), \\
h_2 &= \frac12(s+p_x-p_y-p_z), \\
h_3 &= \frac12(s-p_x+p_y-p_z), \\
h_4 &= \frac12(s-p_x-p_y+p_z).
\end{align*}

もし切断した結合の方向が \(h_1\) なら、フラグメント \(x\) の計算では \(h_1\) を
使わせたくない。このとき、許される方向は \(h_2,h_3,h_4\) である。局所的には
\[
  4\to 3
\]
の次元削減が起きる。

この単純な直交基底の例では、炭素上の変換行列は
\[
\Tmat_{\mathrm C}
=
\begin{array}{c|ccc}
 & h_2 & h_3 & h_4\\
\hline
s   & \frac12 & \frac12  & \frac12 \\
p_x & \frac12 & -\frac12 & -\frac12 \\
p_y & -\frac12 & \frac12 & -\frac12 \\
p_z & -\frac12 & -\frac12 & \frac12
\end{array}
\]
である。

この \(\Tmat_{\mathrm C}\) は、\(h_1\) の列を落として、残りの3本だけを並べたものだと
考えられる。

実際の計算では、基底はもっと大きく、AOも非直交である。そのため、係数はこのような
きれいな \(1/2\) にはならない。現実には \(\Smat\) を使って
\[
  \Amat^\dagger\Smat\Tmat=\mathbf 0
\]
を満たすように \(\Tmat\) を作る。それでも、「切断方向の1本を除き、残りの許された方向
だけでMOを作る」という考え方は同じである。

\section{結合を切ったときの原子核と電子数の扱い}

矩形変換もレベルシフト射影も、操作しているのは電子の変分空間、すなわちMO係数
\(\Cmat\) だけである。\(\alpha\) 炭素の原子核そのものは、これらの操作では一切
分割されない。この点は誤解しやすいので、独立の節で整理する。

\subsection{原子核は分割されない}

\(\alpha\) 炭素で共有結合を切る場合でも、炭素の原子核は点電荷として、どちらか一方の
フラグメントに丸ごと属する。核電荷を半分にしたり、混成軌道の本数に合わせて割ったり
することはない。混成軌道は4本のうち1本(\(h_1\))を除くが、これは電子が動ける方向を
1つ減らすだけであり、核の位置や核電荷 \(Z_{\mathrm C}\) は変わらない。

\begin{verbatim}
   nucleus  Z_C = +6   <- 点電荷。フラグメント x に丸ごと属する(分割しない)
      |
   s, px, py, pz       <- AOは炭素上に4つとも残る
      |
   除くのは h1 という「1方向」だけ(核でもAOそのものでもない)
\end{verbatim}

\(\alpha\) 炭素をフラグメント \(x\) に割り当てるなら、次のように分配される。

\begin{itemize}
\item 核電荷 \(+Z_{\mathrm C}\) はフラグメント \(x\) のコアハミルトニアン
      \(\Hcore\) に入る(核-電子引力として)。
\item 相手フラグメント側の核は、環境静電ポテンシャル \(\Venv\) の側に入る。
\item 炭素上のAO(\(s,p_x,p_y,p_z\) およびその一般化)は、すべてそのまま基底に残る。
      矩形変換で消えるのは、これらのAOが張る空間の中の「1方向」だけである。
\end{itemize}

つまり、核とAOは分割・削除されず、変換されるのは電子が使える方向の集合だけである。

\subsection{AOの本数は減らない。減るのは変分の自由度}

ここは特に混乱しやすいので強調する。矩形変換は、AO基底関数の本数 \(N\) を物理的に
減らすわけではない。基底関数は \(N\) 個のままであり、二電子積分もフルの \(N\) 次元で
計算する。減るのは、SCFがMOを探すときに動かしてよい自由度であり、\(N\) から
\(N-m\) になる。式で言えば、\(\Cmat=\Tmat\Cr\) の \(\Cr\) が \(N-m\) 次元だという
だけである。

\subsection{切った結合の2電子をどちらが持つか}

共有結合を1本切ると、その結合にいた電子対(2電子)の帰属を別途決める必要がある。
混成軌道を除去する操作は「方向」を消すだけで、電子数そのものは自動的には決まらない。
FMOの標準的なやり方では、切断される結合に属する電子対を片方のフラグメントへ割り当て、
境界の炭素(bond-detached atom などと呼ばれる)の扱いをそれに合わせる。

したがって実装では、次の3つを整合させる必要がある。

\begin{enumerate}
\item どちらのフラグメントに \(\alpha\) 炭素の核(と内殻)を割り当てるか。
\item 切断結合の電子対をどちらのフラグメントに割り当てるか。
\item どの混成軌道を変分空間から除くか。
\end{enumerate}

矩形変換 \(\Tmat\)(あるいはレベルシフト \(\Bmat\))が担当するのは3番目だけである。
1番目と2番目は、フラグメント分割の定義として別に決めておく。

\section{\(d\) 軌道の \(6\to5\) 変換との似ている点と違う点}

量子化学では、カーテシアン \(d\) 関数6個から、球面調和 \(d\) 関数5個へ変換することが
ある。これも \(6\to5\) の次元削減なので、形だけ見ると矩形変換に似ている。
ただし、意味は違う。

\begin{center}
\begin{tabular}{|p{0.18\linewidth}|p{0.36\linewidth}|p{0.36\linewidth}|}
\hline
\textbf{比較項目} & \textbf{カーテシアン \(d\) から球面 \(d\)} & \textbf{FMO境界の矩形変換} \\
\hline
次元 & \(6\to5\) & \(N\to N-m\) \\
\hline
落とすもの & \(d\) 関数に混じる冗長な \(s\) 的成分 & 切断結合方向の混成軌道 \\
\hline
行列の決まり方 & 角運動量で決まるほぼ普遍的な変換 & 結合方向、\(\Amat\)、\(\Smat\) に依存 \\
\hline
目的 & 基底関数の表現を整理する & フラグメント境界の不要成分を除去する \\
\hline
\end{tabular}
\end{center}

どちらも「長方形の行列で次元を減らす」という点では似ている。しかし、FMO境界処理の
\(\Tmat\) は、系ごとに決まる物理的な拘束を入れるための行列である。

\section{レベルシフト法との関係}

矩形変換による方法は、レベルシフト法と無関係な別手法ではない。むしろ、次のように
理解できる。
\[
  \boxed{\text{deflation}\ \Longleftrightarrow\ B\to\infty}
\]
つまり、矩形変換はレベルシフト法の \(B\to\infty\) 極限に対応する。

レベルシフト法では、\(B\) を大きくするほど不要な軌道の成分が小さくなる。理想的には
\(B\) を無限大にすれば完全に除去できる。しかし、実際の数値計算で無限大の \(B\) は
使えない。非常に大きな有限値を使うと、条件数が悪くなる。

矩形変換は、この「\(B\) を無限大にした理想状態」を、実際には \(B\) を使わずに実現する。
そのため、「レベルシフト \(B\) をどの値にするか」という調整問題そのものが消える。

\section{行列が「正方でなくなる」のはどこか}

ここで自然に出てくる疑問は、「レベルシフト \(B\) を使う元のやり方でも、行と列の
サイズが変わって正方行列でなくなるのか」というものである。答えを先に言うと、次の
通りである。

\begin{quote}
レベルシフト法では、行列はすべて \(N\times N\) の正方行列のままである。長方形に
なるのは矩形変換のときだけで、しかも長方形になるのは変換行列 \(\Tmat\) 自身であって、
最終的に対角化する行列は再び正方である。
\end{quote}

\subsection{レベルシフト法: 形は \(N\times N\) のまま}

レベルシフト込みの行列
\[
  \widetilde{\Hcore}
  =
  \Hcore+\Venv+\Smat\Amat\Bmat\Amat^\dagger\Smat
\]
は \(N\times N\) の正方行列である。射影項 \(\Smat\Amat\Bmat\Amat^\dagger\Smat\) も、
サイズを追うと
\[
  \underbrace{\Smat}_{N\times N}\,
  \underbrace{\Amat}_{N\times m}\,
  \underbrace{\Bmat}_{m\times m}\,
  \underbrace{\Amat^\dagger}_{m\times N}\,
  \underbrace{\Smat}_{N\times N}
  \;\longrightarrow\; N\times N
\]
となる。途中に \(m\) が現れるので階数(ランク)は \(m\) だが、行数\(\times\)列数という
「形」は \(N\times N\) のままである。

したがって、レベルシフト法では次のようになる。

\begin{itemize}
\item 解く一般化固有値問題は \(N\times N\) のまま。次元は減らない。
\item MOも \(N\) 個出てくる。不要な混成軌道の方向は、エネルギーが \(B\) だけ持ち
      上がった高い仮想軌道として、\(N\) 個の解の中に残る。消えて次元が減るのではなく、
      「上に押しやられて占有されない」状態になるだけである。
\end{itemize}

\subsection{矩形変換: 長方形になるのは \(\Tmat\) だけ}

矩形変換では、変換行列
\[
  \Tmat:\quad N\times(N-m)\qquad(\text{長方形})
\]
が唯一の長方形の行列である。これで挟むと、変換後の行列は
\[
  \Fr=\Tmat^\dagger\Fmat\Tmat:\ (N-m)\times(N-m),\qquad
  \Sr=\Tmat^\dagger\Smat\Tmat:\ (N-m)\times(N-m)
\]
のように、再び正方行列になる。ただしサイズが \(N\) から \(N-m\) へ小さくなっている。

つまり、「行と列のサイズが違う長方形の行列」が現れるのは矩形変換のときの
\(\Tmat\) だけであり、実際に対角化するFock行列・重なり行列は、どちらの方法でも
正方行列である。

\begin{center}
\begin{tabular}{|p{0.26\linewidth}|p{0.3\linewidth}|p{0.3\linewidth}|}
\hline
\textbf{行列} & \textbf{レベルシフト射影} & \textbf{矩形変換によるdeflation} \\
\hline
変換行列 & なし(実質は恒等 \(\mathbf I\), \(N\times N\)) &
\(\Tmat\): \(N\times(N-m)\)(長方形) \\
\hline
対角化する行列 & \(\widetilde{\Hcore}\): \(N\times N\)(正方) &
\(\Fr,\Sr\): \((N-m)\times(N-m)\)(正方) \\
\hline
MOの本数 & \(N\) & \(N-m\) \\
\hline
不要方向の扱い & \(N\) 個の中に高エネルギーで残る & そもそも空間に存在しない \\
\hline
\end{tabular}
\end{center}

\subsection{別の見方: \(\Tmat\) の列数が答えになっている}

\(\Tmat\) が長方形であること自体が、\(m\) 本の方向を落としたという事実の表現である。
正方行列 \(N\times N\) は「全 \(N\) 方向を保つ」変換であり、レベルシフトはこれ
(恒等変換に射影項を足しただけ)に相当する。長方形 \(N\times(N-m)\) は
「\(N-m\) 方向だけを保つ」変換であり、これが deflation の本質である。

前節の対応
\[
  \text{deflation}\ \Longleftrightarrow\ B\to\infty
\]
は、この形の言葉でも言い換えられる。\(B\to\infty\) 極限では、レベルシフトの
\(N\times N\) 空間のうち \(m\) 方向が無限に持ち上がって事実上使えなくなり、残る
\(N-m\) 次元が、矩形変換の縮約空間とちょうど一致する。

\section{両者の比較}

\begin{center}
\begin{tabular}{|p{0.18\linewidth}|p{0.36\linewidth}|p{0.36\linewidth}|}
\hline
\textbf{項目} & \textbf{レベルシフト射影} & \textbf{矩形変換によるdeflation} \\
\hline
基本発想 & 不要方向をエネルギー的に罰する & 不要方向を変数から消す \\
\hline
使う条件 & Fock行列に \(\Smat\Amat\Bmat\Amat^\dagger\Smat\) を足す &
\(\Amat^\dagger\Smat\Tmat=\mathbf 0\) を満たす \(\Tmat\) を使う \\
\hline
\(B\) & 必要 & 不要 \\
\hline
除去の正確さ & 有限 \(B\) では近似的 & 原理的に厳密 \\
\hline
残留成分 & およそ \(O(B^{-1})\) & 機械精度でゼロ \\
\hline
次元 & \(N\) のまま & \(N-m\) \\
\hline
数値安定性 & 大きな \(B\) で条件数が悪化しうる & \(B\) による条件数悪化がない \\
\hline
見方 & \(B\) を大きくして追い出す & \(B\to\infty\) 極限を有限の変換で実現する \\
\hline
\end{tabular}
\end{center}

\section{実装の流れ}

実装するなら、流れは次のようになる。
\begin{enumerate}
\item 切断した結合方向の混成軌道をAO係数として作る: \(\Amat\)。
\item AO重なり行列を用意する: \(\Smat\)。
\item 次の核を張る基底を求める: \(\ker(\Amat^\dagger\Smat)\)。
\item その基底を必要に応じて \(\Smat\)-正規直交化し、
      \(\Tmat^\dagger\Smat\Tmat=\mathbf I\) を満たす \(\Tmat\) を作る。
\item Fock行列や重なり行列を縮約空間へ変換する:
      \[
        \Fr=\Tmat^\dagger\Fmat\Tmat,\qquad
        \Sr=\Tmat^\dagger\Smat\Tmat .
      \]
\item 縮約空間でSCF方程式
      \[
        \Fr\Cr=\Sr\Cr\boldsymbol\varepsilon
      \]
      を解く。
\item 必要に応じて元のAO空間のMO係数や密度に戻す:
      \[
        \Cmat=\Tmat\Cr,\qquad
        \Dmat=\Tmat\Dr\Tmat^\dagger .
      \]
\end{enumerate}
この流れでは、レベルシフト \(B\) はどこにも登場しない。

\section{まとめ}

FMOでは、共有結合を切ってフラグメントを作るため、切断方向の混成軌道をそのフラグメント
のSCFから取り除く必要がある。

従来のレベルシフト法は、不要な混成軌道に大きなエネルギー \(B\) を与えて、被占有軌道
から追い出す。しかし、有限の \(B\) では不要成分が完全には消えず、\(B\) を大きくしすぎると
数値的に不安定になる。

矩形変換では、不要な混成軌道と \(\Smat\)-直交する空間だけを張る行列 \(\Tmat\) を作り、
\[
  \Cmat=\Tmat\Cr
\]
と置く。このとき
\[
  \Amat^\dagger\Smat\Tmat=\mathbf 0
\]
なので、どんな \(\Cr\) を使っても
\[
  \Amat^\dagger\Smat\Cmat=\mathbf 0
\]
が成り立つ。したがって、不要な混成軌道は最初からMOに混ざらない。

この方法は、レベルシフト法の \(B\to\infty\) 極限を、\(B\) を使わずに有限の行列操作で
実現するものである。結果として、レベルシフトの値を調整する必要がなくなり、残留成分と
条件数悪化のトレードオフも避けられる。

\section*{記号一覧}

\begin{longtable}{|p{0.28\linewidth}|p{0.62\linewidth}|}
\hline
\textbf{記号} & \textbf{意味} \\
\hline
\(\Hcore\) & コアハミルトニアン \\
\hline
\(\Venv\) & 他フラグメントなどから来る環境静電ポテンシャル \\
\hline
\(Z_{\mathrm C}\) & 境界(\(\alpha\))炭素の核電荷。分割せず一方のフラグメントに属する \\
\hline
\(\Smat\) & AO重なり行列 \\
\hline
\(\chi_\mu\) & AO基底関数 \\
\hline
\(\psi_p\) & MO \\
\hline
\(\Cmat\) & 元のAO空間でのMO係数 \\
\hline
\(\Cr\) & 縮約空間でのMO係数 \\
\hline
\(\varphi_i^h\) & 除去したい混成軌道 \\
\hline
\(\Amat\) & 除去したい混成軌道のAO係数 \\
\hline
\(m\) & 除去したい混成軌道の本数 \\
\hline
\(N\) & AO基底関数の数 \\
\hline
\(\Bmat\) & レベルシフト行列 \\
\hline
\(\Tmat\) & 許容空間だけを張る矩形変換行列 \\
\hline
\(\Fmat\) & Fock行列 \\
\hline
\(\Gmat[\Dmat]\) & 二電子項 \\
\hline
\(\Dmat\) & 密度行列 \\
\hline
\(\boldsymbol\varepsilon\) & 軌道エネルギー \\
\hline
\end{longtable}

\end{document}
